\documentclass[11pt]{article}
\usepackage[margin=1in]{geometry}
\usepackage{hyperref}
\usepackage{listings}
\usepackage{longtable}
\usepackage{booktabs}
\usepackage{array}
\usepackage{enumitem}
\usepackage{titlesec}
\usepackage{parskip}

\title{FlexQL Server Design Document}
\author{Rahul Das, 25CS60R15 \\ \url{https://github.com/rahul-das-713127/Flex_QL.git}}
\date{\today}

\begin{document}
\maketitle

\tableofcontents
\newpage

\section{Overview}
FlexQL is a lightweight SQL-like database server with a TCP client protocol. The implementation in this repository is a C++17 codebase that focuses on:

\begin{itemize}[leftmargin=*]
\item Persistent storage on disk (table data files + schema catalog, with optional WAL)
\item Fast insert throughput (specialized fast paths, batching, and buffered IO)
\item Basic relational querying (SELECT with single-condition WHERE; INNER JOIN)
\item Performance accelerators (primary-key offset index; small-result caching)
\item Multi-client concurrency (thread-per-connection)
\end{itemize}

This document describes the storage layout, indexing, caching, expiration handling, multithreading, and key design decisions.

\section{Repository Structure}
Project root (\texttt{flexql/}):

\begin{lstlisting}[basicstyle=\ttfamily\small]
flexql/
  Makefile
  server_full.cpp        # Full server implementation (recommended)
  server.cpp             # Simpler/reference server (separate implementation)
  flexql.h               # C client API
  flexql_client.cpp      # Client library implementation
  client_repl.cpp        # REPL client
  benchmark_flexql.cpp   # Benchmark + unit tests client
  server                 # built artifact (make server)
  client                 # built artifact (make client)
  benchmark              # built artifact (make benchmark)
  data/
    catalog.txt          # schema catalog
    wal.log              # write-ahead log (when enabled)
    <TABLE>.data         # per-table data file(s)
\end{lstlisting}

\section{Dependencies and Build Tooling}
\subsection{Language and Compiler}
\begin{itemize}[leftmargin=*]
\item C++17
\item Typical toolchain: \texttt{g++} (or \texttt{clang++})
\end{itemize}

\subsection{System Libraries}
\begin{itemize}[leftmargin=*]
\item POSIX sockets and file APIs (Linux)
\item \texttt{pthread} for threading (linked via \texttt{-lpthread})
\end{itemize}

\subsection{Build System}
\begin{itemize}[leftmargin=*]
\item \texttt{make} via the provided \texttt{Makefile}
\end{itemize}

\section{Data Storage Design}
\subsection{Authoritative Storage is On Disk}
The system does not use RAM as the primary data store. Persistent data resides in the \texttt{data/} directory.

\begin{itemize}[leftmargin=*]
\item \textbf{Schemas}: persisted in \texttt{data/catalog.txt}
\item \textbf{Table rows}: appended to per-table files \texttt{data/<TABLE>.data}
\item \textbf{WAL (optional)}: appended to \texttt{data/wal.log} when enabled
\end{itemize}

RAM is used for speed (buffers, caching, indexes, and metadata), but disk is the source of truth.

\subsection{Schema Catalog (\texttt{catalog.txt})}
Schemas are stored separately from table data. On startup, the server loads schemas from the catalog before serving queries.

Key properties:
\begin{itemize}[leftmargin=*]
\item Schemas are stored under a canonical table name key (case-insensitive handling is implemented by normalizing names).
\item Supported column types include: \texttt{INT}, \texttt{DECIMAL}, \texttt{VARCHAR(n)}, \texttt{DATETIME}.
\item \texttt{EXPIRES\_AT} is mandatory in every table schema.
\end{itemize}

\subsection{Table Data Files (\texttt{<TABLE>.data})}
Each table has a dedicated file created under \texttt{data/}. Inserts are implemented as append-only writes.

\subsubsection{Row Encoding}
Rows are stored in a compact binary format designed for fast append and sequential scan:

\begin{itemize}[leftmargin=*]
\item A header storing \texttt{ncols} (number of columns)
\item For each column, a \texttt{uint32 length} followed by raw bytes of the value
\end{itemize}

Values are stored as their textual form (the same representation used by SQL parsing), which simplifies query execution and makes the storage format schema-agnostic.

\subsubsection{Buffered Writes and Write Arenas}
To reduce syscalls and improve throughput:

\begin{itemize}[leftmargin=*]
\item Inserts serialize into an in-memory per-table \emph{write arena} buffer.
\item Once the arena reaches a threshold, it is flushed to the table \texttt{FILE*}.
\item A large stdio buffer is configured using \texttt{setvbuf(..., \_IOFBF, 8<<20)}.
\item The server uses an async writer thread to flush arena buffers to disk without blocking the hot insert loop.
\end{itemize}

\subsection{Database Recreation vs Persistence}
On startup, the server may wipe \texttt{data/} depending on environment variables.

\begin{itemize}[leftmargin=*]
\item The default and current behavior is to wipe the \texttt{data/} directory on every server start (fresh database per run).
\end{itemize}

This behavior ensures benchmarks and REPL sessions start from a clean database state. The previously supported \texttt{FLEXQL\_PRESERVE\_DATA} environment variable has been removed in favor of this always-fresh semantic.

\section{Write-Ahead Log (WAL)}
\subsection{Purpose}
When enabled, WAL provides crash-recovery by logging SQL statements so that the server can replay them on startup.

\textbf{Important scope:}
\begin{itemize}[leftmargin=*]
\item WAL improves durability and crash recovery.
\item WAL does \emph{not} provide replication, failover, or protection from disk loss.
\end{itemize}

\subsection{Enable/Disable}
\begin{itemize}[leftmargin=*]
\item Default: WAL disabled.
\item Enable: \texttt{FLEXQL\_ENABLE\_WAL=1}
\item Force disable: \texttt{FLEXQL\_DISABLE\_WAL=1}
\end{itemize}

\subsection{Record Format}
The WAL file (\texttt{data/wal.log}) is a binary stream of records:

\begin{longtable}{@{}ll@{}}
\toprule
Field & Description\\
\midrule
\texttt{uint32 len} & Number of bytes in the SQL statement\\
\texttt{len bytes sql} & SQL statement bytes (typically includes trailing semicolon)\\
\texttt{uint32 crc} & CRC32 of the SQL bytes\\
\bottomrule
\end{longtable}

\subsection{WAL Buffers and Flushing}
To minimize lock contention and syscall overhead:

\begin{itemize}[leftmargin=*]
\item Each client handler thread accumulates WAL records into a per-connection buffer.
\item Periodically, the buffer is appended to a global WAL batch buffer under the global mutex.
\item The batch is flushed when size/count thresholds are met, and also on disconnect.
\end{itemize}

\subsection{Replay on Startup}
On startup, the server replays \texttt{wal.log}:

\begin{itemize}[leftmargin=*]
\item The server reads records sequentially.
\item CRC is verified; on corruption/torn tail it stops and truncates the WAL to the last good record.
\item Each SQL statement is executed via the normal execution engine.
\item During replay, a flag disables WAL re-appending to prevent infinite log growth.
\end{itemize}

\textbf{Operational note:} For WAL replay to restore data, \texttt{data/} must be preserved across restarts (\texttt{FLEXQL\_PRESERVE\_DATA=1}). If the server wipes \texttt{data/} at startup, there is nothing to replay.

\section{SQL Layer and Schema Enforcement}
\subsection{CREATE TABLE}
CREATE TABLE parses and stores the schema. The server enforces:

\begin{itemize}[leftmargin=*]
\item Only supported types are allowed: INT, DECIMAL, VARCHAR(n), DATETIME.
\item VARCHAR requires \texttt{n > 0}.
\item \texttt{EXPIRES\_AT} must exist.
\end{itemize}

\subsection{INSERT}
INSERT supports inserting rows into a table. The implementation enforces:

\begin{itemize}[leftmargin=*]
\item Column count matches the schema.
\item Each column value matches its declared type.
\item VARCHAR length is limited to \texttt{n}.
\item \texttt{EXPIRES\_AT} must be a valid positive integer timestamp.
\end{itemize}

Validation is applied for both:
\begin{itemize}[leftmargin=*]
\item Normal insert path
\item Fast insert path
\item Ultra-fast bulk insert path
\end{itemize}

The server supports \textbf{multi-row inserts} of the form:
\texttt{INSERT INTO T VALUES (..),(..),(..);}
which allows client-side batching to reduce round trips.

\subsection{SELECT}
SELECT supports:

\begin{itemize}[leftmargin=*]
\item \texttt{SELECT * FROM T;}
\item \texttt{SELECT C1, C2 FROM T;}
\item Optional single-condition WHERE clause
\end{itemize}

\subsection{WHERE and JOIN Operators}
The condition operator set supported in WHERE and JOIN includes:

\begin{itemize}[leftmargin=*]
\item \texttt{=, >, <, >=, <=}
\end{itemize}

INNER JOIN is supported with an ON condition and an optional WHERE clause.

\section{Expiration Timestamp Handling}
\subsection{EXPIRES\_AT Column}
Every table must contain an \texttt{EXPIRES\_AT} column. Inserts require a valid expiration value.

\subsection{Query-Time Filtering}
During scans (including SELECT and JOIN), expired rows are filtered out. In other words:

\begin{itemize}[leftmargin=*]
\item Rows with \texttt{EXPIRES\_AT} in the past are skipped.
\item Expiration enforcement happens at read time.
\end{itemize}

\subsection{Design Decision}
The server does not run a background vacuum to delete expired rows from disk. Instead:

\begin{itemize}[leftmargin=*]
\item Reads ignore expired rows.
\item Storage usage can grow unless compaction/vacuum is added.
\end{itemize}

\section{Indexing}
\subsection{Primary Key Offset Index}
To accelerate equality lookups on a likely primary key column:

\begin{itemize}[leftmargin=*]
\item The server detects a ``probable PK'' column (typically the first column or a column named \texttt{ID}).
\item If it is INT, the server enables an offset index: \texttt{pk\_offsets}.
\item The index maps PK value $\rightarrow$ file offset for fast retrieval.
\end{itemize}

\subsection{Index Maintenance}
\begin{itemize}[leftmargin=*]
\item Inserts update index state (or mark it invalid for deferred rebuild depending on the path).
\item On startup (after loading schemas), the server can rebuild indexes by scanning table files.
\end{itemize}

\subsection{Index Usage in Queries}
\item \textbf{Index Usage in Queries}: If a WHERE clause matches the PK column with \texttt{=}, \texttt{>}, \texttt{<}, \texttt{>=}, or \texttt{<=}, the query executor uses the \texttt{pk\_offsets} array to skip full-table scans.
\item \textbf{JOIN Acceleration}: INNER JOIN also uses the PK index of the right-hand table for \texttt{=} conditions to perform fast lookups.

\section{Caching Strategy}
\subsection{Goal}
Reduce latency for repeated SELECT queries returning small result sets.

\subsection{Mechanism}
\begin{itemize}[leftmargin=*]
\item An in-memory cache stores results keyed by the exact SQL string.
\item Cache is an LRU-like structure with bounded capacity.
\item Cache entries are used only for small result sets to prevent excessive memory usage.
\end{itemize}

\subsection{Invalidation}
Writes invalidate cached results:

\begin{itemize}[leftmargin=*]
\item Any INSERT operation sets a global dirty flag (atomic).
\item SELECT queries check this flag and clear the cache if it is set.
\end{itemize}

\section{Multithreading and Concurrency}
\subsection{Threading Model}
\begin{itemize}[leftmargin=*]
\item The server listens on TCP port 9000.
\item Each accepted connection is handled by a dedicated \texttt{std::thread} running \texttt{handle\_client}.
\end{itemize}

\subsection{Synchronization}
\begin{itemize}[leftmargin=*]
\item \textbf{Global mutex} protects shared metadata (schemas, runtime maps, WAL batch buffer).
\item \textbf{Per-table spinlock} (using \texttt{std::atomic\_flag}) is used in the hot insert path to reduce overhead under heavy insert loads.
\end{itemize}

\subsection{Write Path Concurrency}
\begin{itemize}[leftmargin=*]
\item For repeated inserts into the same table, the handler can hold the per-table spinlock across multiple rows received in a burst.
\item The write arena reduces file IO contention and amortizes syscall overhead.
\end{itemize}

\subsection{Read Path Concurrency}
\begin{itemize}[leftmargin=*]
\item SELECT execution requires flushing pending write arenas before reads for visibility.
\item This is coordinated via flush hooks and the global mutex.
\end{itemize}

\section{Operational Guide}
\subsection{Environment Variables}
\begin{longtable}{@{}lp{0.72\textwidth}@{}}
\toprule
Variable & Meaning \\
\midrule
\texttt{FLEXQL\_ENABLE\_WAL=1} & Enable WAL logging to \texttt{data/wal.log}. \\
\texttt{FLEXQL\_DISABLE\_WAL=1} & Force WAL disabled. \\
\texttt{FLEXQL\_PIPELINE\_INSERTS=1} & Client-side optimization: pipeline inserts. \\
\texttt{FLEXQL\_BULK\_ACK\_INSERTS=1} & Client/server optimization: bulk acknowledgements. \\
\texttt{FLEXQL\_INSERT\_BATCH\_SIZE} & Benchmark client batch size (default 1). \\
\bottomrule
\end{longtable}

\subsection{Stopping a Server Running in Background}
If the server is running in the background and you see errors like \texttt{bind failed}, find and kill the process:

\paragraph{Using ss}
\begin{lstlisting}[basicstyle=\ttfamily\small]
ss -ltnp | grep ":9000"
# identify PID, then:
kill <PID>
# if it won't stop:
kill -9 <PID>
\end{lstlisting}

\paragraph{Using lsof}
\begin{lstlisting}[basicstyle=\ttfamily\small]
lsof -iTCP:9000 -sTCP:LISTEN -n -P
kill <PID>
\end{lstlisting}

\section{Known Limitations and Future Improvements}
\begin{itemize}[leftmargin=*]
\item Expired rows are filtered at read time; there is no background vacuum/compaction.
\item WAL replay is SQL-replay based; it re-executes statements rather than replaying page-level changes.
\item Indexing is currently focused on a probable INT primary key column and limited to 20M rows.
\item Query planner is minimal (single WHERE condition; no AND/OR).
\end{itemize}

\end{document}
